\section{Microservice implementation: Ingestion and Control Module}
\label{sec:ingestion-control}

Both are written in Java 17 with Javalin 6.1.3 and Gson, like the rest, and share \code{common}: the HTTP server
with its error format (\code{ServiceServer}), the configuration from the environment (\code{Env}), the range datalake
and the coordination library. Ingestion implements \code{api/ingestion-openapi.yaml} and the Control
\code{api/control-openapi.yaml}; both are checked against their specification by tests (\code{OpenApiContractTest})
and by \code{scripts/validate-api.sh} on the running stack.

\subsection{Ingestion Service}

\begin{table}[H]
\centering\small
\begin{tabularx}{\linewidth}{@{}L{4.1cm}X@{}}
\toprule
\textbf{Class} & \textbf{Responsibility} \\
\midrule
\code{IngestionApi} & \code{POST /ingest/\{id\}}, \code{GET /ingest/status/\{id\}}, \code{GET /ingest/list}, \code{GET /status} \\
\code{IngestionService} & One ingestion, without HTTP: locate, fetch, split, save. The API and the workers run this same code \\
\code{GutenbergClient}, \code{LocalFileSource} & The two sources (\code{INGESTION\_SOURCE}): the download, or \code{pg<id>.txt} files of a folder \\
\code{BookSplitter} & Header and body between the \code{*** START} and \code{*** END} markers, migrated from Stage 1 \\
\code{IngestionWorkers}, \code{IngestionWorker} & The \code{INGESTION\_WORKERS} consumers of \code{ingest-queue} \\
\code{Throttle} & One pause for the whole replica after a 429 or a 403 of Gutenberg \\
\code{BackgroundConnection} (common) & Connects to Hazelcast without blocking the start, and follows the connection \\
\bottomrule
\end{tabularx}
\caption{Structure of the Ingestion Service.}
\end{table}

\textbf{Two ways in, one code path.} \code{POST /ingest/\{id\}} is synchronous: when it answers 200 the book is
complete in the datalake (\code{downloaded}, or \code{already\_available} if it was there). It is for administration
and for the validation. In normal operation the books arrive through \code{ingest-queue}, and each worker calls the
same \code{IngestionService}, so both paths store a book in exactly the same way.

\textbf{The source.} \code{GutenbergClient} downloads \url{https://www.gutenberg.org/cache/epub/<id>/pg<id>.txt}
with an identifying User-Agent. Its timeout (30~s) bounds the whole download: the timeout of the JDK client alone
covers only the headers, and a body that arrives slowly was measured to last far longer. A 404 is
\code{NOT\_FOUND}, a permanent error; a timeout, a network error or a 5xx is a transient one; a 429 or a 403 means that
Gutenberg is throttling us. \code{LocalFileSource} reads the same text from a folder, untouched, so the benchmarks and
the synthetic dataset (DR-16) run every step after the fetch exactly as online, without the network.

\subsubsection{Readiness of the datalake}
\label{sec:readiness}

A book is in the datalake only when its \code{<id>.\_READY} marker exists together with its header and its body
(\code{api/contracts/datalake-readiness.md}). Each file is written to a temporary file with a unique name and moved
into place with an atomic move, and the marker is written last, with the time of the save. So:
\begin{itemize}
  \item a process that dies halfway leaves files without a marker: the book does not exist for anyone, and the next
    ingestion writes it again, whole;
  \item Indexing checks the marker before reading a book (\code{NOT\_READY} otherwise), and the Control's recovery
    counts only books with a marker;
  \item a ready book is never written again: \code{save} is a no-op and the ingestion answers \code{already\_available}
    without fetching anything.
\end{itemize}
\code{locate} computes the path from the id, without listing any folder, so checking one book costs the same with 200
books or with 100,000.

\subsubsection{Task ownership and leases}
\label{sec:ownership}

A worker owns a book from its claim to its completion. The claim polls \code{ingest-queue} and writes
\code{INGESTING} with the worker's name (\code{ingestion-<hostname>/w<n>}; inside Docker the hostname is the container
id) and \code{leaseUntil = now + LEASE\_SECONDS} (120~s), in one transaction. If the process dies after the poll, the
transaction is rolled back and the id goes back to the queue; if it dies after the commit, the book stays
\code{INGESTING} with a lease that nobody renews.

The Control queries \code{state IN (INGESTING, INDEXING) AND leaseUntil < now} every \code{LEASE\_CHECK\_SECONDS}
(10~s) and requeues those books with one more attempt (\code{LEASE\_EXPIRED}). The completion is a compare-and-set
against the entry the worker claimed: if the Control requeued the book meanwhile, the worker's completion fails, it
drops the book, and whatever it wrote is harmless because writing is idempotent. The library can renew a lease
(\code{Coordination.renew}), but Ingestion does not need to: one attempt lasts at most $3 \times 30 + 3$~s with the defaults (the download and two retries, each bounded by the
30~s timeout, plus their pauses of 1 and 2~s), well within the 120~s lease. So a lease only expires when a process
really died or hung.

\subsubsection{Retries and errors}

Each failure goes back to the contract with a code (§4.3 of the coordination contract):
\begin{itemize}
  \item \textbf{Permanent} (\code{NOT\_FOUND}, \code{INVALID\_BOOK} for a text without markers): \code{FAILED} at once.
    Retrying would give the same answer.
  \item \textbf{Transient} (\code{DOWNLOAD\_FAILED}): first \code{INGESTION\_FETCH\_RETRIES} (2) retries inside the
    attempt, after 1~s and 2~s; then the book goes back to the queue with one more attempt, and after
    \code{MAX\_ATTEMPTS} (3) it is \code{FAILED}. \code{POST /control/retry-failed} gives the failed books another round.
  \item \textbf{Throttling} (429 or 403): not retried. The book goes back to the queue as a transient failure, and the
    whole replica stops claiming for a pause of 30~s that doubles at each new throttling, up to 10~min (a successful
    download brings it back to 30~s). Otherwise every worker would keep hitting Gutenberg and burn the attempts of every
    book in the queue in a few seconds; with the pause they keep them.
  \item \textbf{Hazelcast down:} the worker waits and tries again with a growing pause (1 to 30~s); the client
    reconnects by itself.
\end{itemize}

\subsubsection{Idempotency}

Repeating any step leaves the same result: seeding uses \code{putIfAbsent}; a book already ready is not fetched; each
file is replaced atomically by a complete version, so two ingestions of the same book at the same time (an
administrator's \code{POST /ingest} and a worker) leave one complete copy; and an id that appears twice in the queue is
dropped by the claim, because its entry is no longer \code{PENDING} (counted in \code{ids\_skipped}). This is what makes
the retries and the lease recovery safe: a book processed twice costs time, never a wrong result.

\textbf{Start and status.} The HTTP server starts first, then the Hazelcast client connects in the background
(\code{BackgroundConnection}, shared with the Control and Indexing), and the workers start once it is connected.
\code{GET /status} answers from the first moment and shows \code{coordination} (\code{connecting},
\code{connected}, \code{disconnected}, \code{failed} or \code{closed}) and the counters of the workers, including
\code{coordination\_cost}: the wait in the queue, the claim, the work and the completion of each book (CF-12). A
\code{SIGTERM} (\code{docker stop}) lets each worker finish its book for up to 10~s before the client and the server
close.

\subsection{Control Module}

\begin{table}[H]
\centering\small
\begin{tabularx}{\linewidth}{@{}L{3.6cm}X@{}}
\toprule
\textbf{Class} & \textbf{Responsibility} \\
\midrule
\code{ControlApi} & \code{POST /control/jobs}, \code{GET /control/status}, \code{GET /control/tasks/\{id\}},
  \code{POST /control/retry-failed} \\
\code{ControlTasks} & The lease check and the decision to publish \\
\code{ControlScheduler} & Runs both tasks on their own threads; a failure of one run is logged and the next run goes on \\
\code{Recovery}, \code{StorageEvidence} & Rebuild \code{book-state} when the cluster lost it (CF-10) \\
\code{IndexingClient} & \code{POST /index/publish}, with a 30~s timeout \\
\bottomrule
\end{tabularx}
\caption{Structure of the Control Module.}
\end{table}

\textbf{The API.} \code{POST /control/jobs} takes 1 to 100,000 ids and answers 202 with how many were new and how many
were already known: it only seeds, the work happens asynchronously in the workers. \code{GET /control/status} counts
the books per state and the ids in each queue, and shows the last publication, the last lease check and the last
recovery. \code{GET /control/tasks/\{id\}} shows the entry of one book: owner, lease, attempts and last error. While the
Control is not connected to Hazelcast these routes answer 503 \code{COORDINATION\_UNAVAILABLE} at once instead of
waiting for the cluster.

\textbf{Publication of completed index updates.} The Control decides \emph{when} to publish; Indexing does the
publication. Every second the Control counts the \code{INDEXED} books, and calls \code{POST /index/publish} when there are
\code{PUBLISH\_BATCH} of them (100) or when \code{PUBLISH\_INTERVAL\_S} (30~s) have passed with at least one. A failed call
(Indexing down, a timeout, 409 during a rebuild, 503 without MongoDB) is retried after 1, 2, 4\ldots{} up to 30~s, not
every second; the publication is idempotent, so a retry never duplicates a version. A publication that makes nothing
new visible (the books were already published but T10 has not marked them yet) does not trigger another one until the
interval is over, or until the \code{INDEXED} books drop below the batch.

\subsubsection{Recovery when Hazelcast loses its state}
\label{sec:recovery}

Hazelcast keeps everything in memory, so if both members stop the state of every book is gone. The Control writes a
marker in \code{coordination-meta} after its first seed or recovery; at start and before every lease check it looks
for it. If it is missing, it rebuilds \code{book-state} from the persistent evidence, from the most advanced state
down (contract §7, row R):
\begin{enumerate}
  \item the books of \code{index\_books} with a version up to \code{CURRENT} are \code{PUBLISHED};
  \item the rest of \code{index\_log} are \code{INDEXED};
  \item the rest with a \code{\_READY} marker are \code{DOWNLOADED}, and offered again to \code{index-queue};
  \item the seed file, if any, is seeded again as \code{PENDING};
  \item the marker is written.
\end{enumerate}
Every source is read before anything is written, a recovered book never goes backwards, and an entry that already
exists is only moved forward from \code{PENDING}. So a recovery interrupted halfway can simply run again. What has no
evidence is lost from the cluster: a book that was \code{PENDING}, or \code{INGESTING} without its marker. It is sent again
with \code{POST /control/jobs}, which leaves the recovered books as they are. The failure tests of CF-11 restart the whole
cluster in the middle of a run and check this end to end (Section~\ref{sec:failure-tests}).
